\ifx\type\undefined
  \documentclass[10pt, t]{beamer}
  \setbeamertemplate{footline}[page number]
\else
  \documentclass[10pt]{article}
  \usepackage[margin=1in]{geometry}
\fi

\usepackage{amsmath}
\usepackage{amssymb}
\usepackage{amsthm}
\usepackage{cancel}
\usepackage{listings}
\usepackage{mathrsfs}
\usepackage{multirow}
\usepackage{soul}
\usepackage{stmaryrd}
\usepackage{tikz}
\usepackage{tikz-cd}
\usepackage{wrapfig}

\newtheorem*{algorithm}{Algorithm}
\newtheorem*{assumptions}{Assumptions}
\newtheorem*{conjecture}{Conjecture}
\newtheorem*{consequences}{Consequences}
\newtheorem*{exercise}{Exercise}
\newtheorem*{formalisation}{Formalisation}
\newtheorem*{proposition}{Proposition}
\newtheorem*{question}{Question}
\newtheorem*{remark}{Remark}

\ifx\type\undefined\else
  \newtheorem*{definition}{Definition}
  \newtheorem*{example}{Example}
  \newtheorem*{lemma}{Lemma}
  \newtheorem*{theorem}{Theorem}
\fi

\definecolor{keywordcolor}{rgb}{0.7, 0.1, 0.1}
\definecolor{tacticcolor}{rgb}{0.0, 0.1, 0.6}
\definecolor{commentcolor}{rgb}{0.4, 0.4, 0.4}
\definecolor{symbolcolor}{rgb}{0.0, 0.1, 0.6}
\definecolor{sortcolor}{rgb}{0.1, 0.5, 0.1}
\definecolor{attributecolor}{rgb}{0.7, 0.1, 0.1}
\def\lstlanguagefiles{lstlean.tex}
\lstset{language=lean}

\newcommand\A{\mathbb{A}}
\newcommand\C{\mathbb{C}}
\newcommand\F{\mathbb{F}}
\newcommand\G{\mathbb{G}}
\renewcommand\H{\mathbb{H}}
\newcommand\I{\mathbb{I}}
\newcommand\N{\mathbb{N}}
\renewcommand\P{\mathbb{P}}
\newcommand\Q{\mathbb{Q}}
\newcommand\R{\mathbb{R}}
\newcommand\Z{\mathbb{Z}}

\renewcommand\AA{\mathcal{A}}
\newcommand\BB{\mathcal{B}}
\newcommand\CC{\mathcal{C}}
\newcommand\DD{\mathcal{D}}
\newcommand\EE{\mathcal{E}}
\newcommand\FF{\mathcal{F}}
\newcommand\GG{\mathcal{G}}
\newcommand\HH{\mathcal{H}}
\newcommand\II{\mathcal{I}}
\newcommand\LL{\mathcal{L}}
\newcommand\MM{\mathcal{M}}
\newcommand\NN{\mathcal{N}}
\newcommand\OO{\mathcal{O}}
\newcommand\PP{\mathcal{P}}
\newcommand\RR{\mathcal{R}}
\renewcommand\SS{\mathcal{S}}
\newcommand\TT{\mathcal{T}}
\newcommand\XX{\mathcal{X}}

\renewcommand\aa{\mathfrak{a}}
\newcommand\cc{\mathfrak{c}}
\newcommand\dd{\mathfrak{d}}
\newcommand\ff{\mathfrak{f}}
\renewcommand\gg{\mathfrak{g}}
\newcommand\mm{\mathfrak{m}}
\newcommand\pp{\mathfrak{p}}
\newcommand\qq{\mathfrak{q}}
\renewcommand\ss{\mathfrak{s}}

\newcommand\LLL{\mathscr{L}}

\newcommand\ab{\mathrm{ab}}
\newcommand\Ab{\mathbf{Ab}}
\newcommand\Alg{\mathbf{Alg}}
\newcommand\Aff{\mathbf{Aff}}
\newcommand\Aut{\operatorname{Aut}}
\newcommand\Az{\mathrm{Az}}
\newcommand\Br{\operatorname{Br}}
\newcommand\BSD{\operatorname{BSD}}
\newcommand\ch{\operatorname{char}}
\newcommand\Cl{\operatorname{Cl}}
\newcommand\coker{\operatorname{coker}}
\newcommand\contr{\operatorname{contr}}
\newcommand\cris{\mathrm{cris}}
\renewcommand\d{\mathrm{d}}
\newcommand\Div{\operatorname{Div}}
\newcommand\dR{\mathrm{dR}}
\newcommand\EN{\operatorname{EN}}
\newcommand\End{\operatorname{End}}
\newcommand\ES{\operatorname{ES}}
\newcommand\et{\mathrm{\acute{e}t}}
\newcommand\Et{\mathbf{\acute{E}t}}
\newcommand\Ext{\operatorname{Ext}}
\newcommand\Fr{\operatorname{Fr}}
\newcommand\Frac{\operatorname{Frac}}
\newcommand\Gal{\operatorname{Gal}}
\newcommand\GL{\operatorname{GL}}
\newcommand\Gr{\mathrm{Gr}}
\newcommand\Hom{\operatorname{Hom}}
\newcommand\HT{\mathrm{HT}}
\newcommand\id{\operatorname{id}}
\newcommand\im{\operatorname{im}}
\newcommand\Ind{\operatorname{Ind}}
\renewcommand\inf{\operatorname{inf}}
\newcommand\inv{\operatorname{inv}}
\newcommand\Irr{\operatorname{Irr}}
\newcommand\Jac{\operatorname{Jac}}
\newcommand\lcm{\operatorname{lcm}}
\newcommand\Mat{\operatorname{Mat}}
\newcommand\Mod{\mathbf{Mod}}
\newcommand\Nm{\operatorname{Nm}}
\newcommand\nr{\mathrm{nr}}
\newcommand\NS{\operatorname{NS}}
\newcommand\Ob{\operatorname{Ob}}
\newcommand\ord{\operatorname{ord}}
\newcommand\op{\mathrm{op}}
\newcommand\PGL{\operatorname{PGL}}
\newcommand\Pic{\operatorname{Pic}}
\newcommand\Prob{\operatorname{Prob}}
\newcommand\Proj{\operatorname{Proj}}
\newcommand\PSh{\mathbf{PSh}}
\newcommand\Reg{\operatorname{Reg}}
\newcommand\res{\operatorname{res}}
\newcommand\rk{\operatorname{rk}}
\newcommand\Sch{\mathbf{Sch}}
\newcommand\Sel{\operatorname{Sel}}
\newcommand\Set{\mathbf{Set}}
\newcommand\sgn{\operatorname{sgn}}
\newcommand\Sh{\mathbf{Sh}}
\newcommand\SL{\operatorname{SL}}
\newcommand\Spec{\operatorname{Spec}}
\newcommand\supp{\operatorname{supp}}
\newcommand\Tam{\operatorname{Tam}}
\newcommand\Top{\mathbf{Top}}
\newcommand\tor{\operatorname{tor}}
\newcommand\tr{\operatorname{tr}}
\newcommand\tra{\operatorname{tra}}
\newcommand\WC{\operatorname{WC}}

\DeclareFontFamily{U}{wncyr}{}
\DeclareFontShape{U}{wncyr}{m}{n}{<->wncyr10}{}
\DeclareSymbolFont{cyr}{U}{wncyr}{m}{n}
\DeclareMathSymbol{\Sha}{\mathord}{cyr}{"58}

\newcommand{\function}[5][]{
  \if &#1&
    \begin{array}{rcl}
      #2 & \longrightarrow & #3 \\
      #4 & \longmapsto     & #5
    \end{array}
  \else
    \begin{array}{rcrcl}
      #1 & : & #2 & \longrightarrow & #3 \\
         &   & #4 & \longmapsto     & #5
    \end{array}
  \fi
}

\newcommand{\functions}[7][]{
  \if &#1&
    \begin{array}{rcl}
      #2 & \longrightarrow & #3 \\
      #4 & \longmapsto     & #5 \\
      #6 & \longmapsto     & #7 \\
    \end{array}
  \else
    \begin{array}{rcrcl}
      #1 & : & #2 & \longrightarrow & #3 \\
         &   & #4 & \longmapsto     & #5 \\
         &   & #6 & \longmapsto     & #7
    \end{array}
  \fi
}
\usetheme{szeged}
\title{Torsion and Tamagawa numbers over function fields}
\subtitle{Modular Curves and Applications of AI to Number Theory}
\author{David Kurniadi Angdinata (with Mentzelos Melistas)}
\institute{University of East Anglia}
\date{Thursday, 17 September 2026}

\begin{document}

\frame{\titlepage}

\section{Number fields}

\begin{frame}{The Birch and Swinnerton-Dyer conjecture}

Let $ E $ be an elliptic curve over a global field $ K $.

\bigskip Tate 1966 conjectured that $ \ord_{s = 1} L(E, s) = \rk(E) $, and that
$$ L^*(E) = \dfrac{\Omega(E) \cdot \Reg(E) \cdot \#\Sha(E) \cdot c(E)}{\mu(\A_K / K) \cdot \#\tor(E)^2}. $$
Here, $ \tor(E) $ is the \textbf{torsion subgroup}, and the \textbf{Tamagawa number} is
$$ c(E) := \prod_v c_v(E) := \prod_v [E(K_v) : E_0(K_v)], $$
where $ c_v(E) $ counts the $ k_v $-rational components of its N\'eron model at $ v $.

\bigskip Are there any cancellations between $ \#\tor(E)^2 $ and $ c(E) $?

\end{frame}

\begin{frame}{Known divisibilities}

Let $ E $ be an elliptic curve over a number field $ K $.

\begin{theorem}[Lorenzini 2011]
If $ K = \Q $ and $ \tor(E) \cong T $, then $ c $ divides $ c(E) $.
\begin{center}
\begin{tabular}{|c|c|l|}
\hline
$ T $ & $ c $ & exceptions \\
\hline
$ C_4 $, $ C_2 \oplus C_4 $ & $ 2 $ & $ X_1(15) $, 15a7, 17a4 \\
\hline
$ C_6 $, $ C_2 \oplus C_6 $ & $ 6 $ & $ X_1(14) $, 14a6, 20a2 \\
\hline
$ C_5 $, $ C_7 $, $ C_8 $, $ C_9 $, & \multirow{2}{*}{$ \#T $} & \multirow{2}{*}{$ X_1(11) $} \\
$ C_{10} $, $ C_{12} $, $ C_2 \oplus C_8 $ & & \\
\hline
\end{tabular}
\end{center}
If $ K = \Q(\sqrt{d}) $ and $ E $ has a point of order $ N = 11, 13 $, then $ N $ divides $ c(E) $. In general, if $ E $ has a point of order $ N = 7, 9 $, then $ N $ divides $ c(E) $, up to at most $ 2^{9(\rk\OO_K^\times + 1)} $ exceptions.
\end{theorem}

\end{frame}

\begin{frame}{Proof strategy}

Let $ P $ be a point of order $ N \ge 4 $. Then $ E $ has a \textbf{Tate normal form}
$$ \TT(a, b) : y^2 + (1 - a)xy - by = x^3 - bx^2, \qquad a, b \in K, $$
with $ P = (0, 0) $, and $ [N](0, 0) = O $ gives an affine model of $ X_1(N) $.

\bigskip If $ v(a), v(b) > 0 $ for some place $ v $ of $ K $, then the reduction of $ E $ at $ v $ is $ y^2 + xy = x^3 $, which is split multiplicative with $ c_v(E) = v(\Delta) $.

\bigskip

\begin{example}[$ K = \Q $ and $ N = 5 $]
The condition $ 2P = -3P $ forces $ a = b =: f $ with $ \Delta = f^5(f^2 - 11f - 1) $. If $ v(f) \ne 0 $ for some place $ v $ of $ \Q $, then assuming $ v(f) > 0 $ gives $ c_v(E) = 5v(f) $. Otherwise, $ f \in \Z^\times = \{\pm1\} $ forces $ E = X_1(11) $.
\end{example}

\end{frame}

\section{Global function fields}

\begin{frame}{Known divisibilities}

Let $ K $ be the function field of a nice curve $ C $ of genus $ g $ over $ k = \F_q $, and let $ E $ be a non-isotrivial elliptic curve over $ K $ of height $ h = \deg\omega_E > 0 $.

\bigskip Then $ \#\tor(E) $ divides $ c(E) $ easily, by considering the map
$$ \tor(E) \hookrightarrow E(K) \hookrightarrow \prod_v E(K_v) \to \prod_v E(K_v) / E_0(K_v), $$
which is injective by Shioda's height formula
$$ \langle P, P\rangle = 2h + 2(P \cdot O) - \sum_v \contr_v(P), \qquad P \in E(K). $$
Ulmer 2014 also proved that $ \#\tor(E)^2 $ divides $ \Reg(E) \cdot D(E) $, where $ D(E) $ is the absolute discriminant of the trivial lattice over $ k $ of the minimal proper regular model of $ E $, which is related to $ c(E) $.

\end{frame}

\begin{frame}{New divisibilities}

\begin{theorem}[A.--Melistas 2026]
If $ K = k(t) $ with $ \ch(k) \ne 2, 3 $, then $ \#\tor(E)^2 $ divides $ c(E) $, except when $ \#\tor(E) = 2, 3, 4 $.
\end{theorem}

\bigskip Since $ \Omega(E) = q^{-h} $ and $ \mu(\A_K / K) = q^{-1} $,
$$ \dfrac{q^h}{\Reg(E)} \cdot L^*(E) \overset{\BSD}{=} \dfrac{q \cdot \#\Sha(E) \cdot c(E)}{\#\tor(E)^2} \in \Z[\tfrac{1}{6}]. $$

\begin{theorem}[A.--Melistas 2026]
A point of prime order $ 7 \le N \le 101 $ with $ N \nmid q $ forces $ N^2 $ to divide $ c(E) $.
\end{theorem}

\bigskip This should hold for all primes $ N \ge 5 $, including $ N \mid q $.

\end{frame}

\begin{frame}{Sharp divisibilities}

\begin{center}
\scriptsize
\renewcommand{\arraystretch}{0.88}
\begin{tabular}{|c|c|c|c|c|c|c|}
\cline{1-3}\cline{5-7}
$ T $ & $ \ch(k) $ & $ c $ & & $ T $ & $ \ch(k) $ & $ c $ \\
\cline{1-3}\cline{5-7}
$ C_2 $ & any & $ 2 $ & & $ C_2 \oplus C_2 $ & $ \ne 2 $ & $ 2^3 $ \\
\cline{1-3}\cline{5-7}
$ C_3 $ & any & $ 3 $ & & $ C_3 \oplus C_3 $ & $ \ne 3 $ & $ 3^4 $ \\
\cline{1-3}\cline{5-7}
\multirow{2}{*}{$ C_4 $} & $ \ne 2 $ & $ 2^3 $ & & $ C_4 \oplus C_2 $ & $ \ne 2 $ & $ 4^2 \cdot 2^2 $ \\
\cline{2-3}\cline{5-7}
& $ 2 $ & $ 2^2 $ & & \multirow{2}{*}{$ C_6 \oplus C_2 $} & $ \ne 2, 3 $ & $ 2^3 \cdot 6^3 $ \\
\cline{1-3}\cline{6-7}
$ C_5 $ & any & $ 5^2 $ & & & $ 3 $ & $ 6^3 $ \\
\cline{1-3}\cline{5-7}
\multirow{3}{*}{$ C_6 $} & $ \ne 2, 3 $ & $ 6^2 $ & & $ C_8 \oplus C_2 $ & $ \ne 2 $ & $ 2^{16} $ \\
\cline{2-3}\cline{5-7}
& $ 2 $ & $ 2 \cdot 6 $ & & $ C_4 \oplus C_4 $ & $ \ne 2 $ & $ 4^6 $ \\
\cline{2-3}\cline{5-7}
& $ 3 $ & $ 3 \cdot 6 $ & & \multirow{2}{*}{$ C_6 \oplus C_3 $} & $ \ne 2, 3 $ & $ 3^4 \cdot 6^4 $ \\
\cline{1-3}\cline{6-7}
$ C_7 $ & any & $ 7^3 $ & & & $ 2 $ & $ 6^4 $ \\
\cline{1-3}\cline{5-7}
\multirow{2}{*}{$ C_8 $} & $ \ne 2 $ & $ 8^3 $ & & $ C_5 \oplus C_5 $ & $ \ne 5 $ & $ 5^{12} $ \\
\cline{2-3}\cline{5-7}
& $ 2 $ & $ 8^2 $ & & $ C_{11} $ & $ 11 $ & $ 11^5 $ \\
\cline{1-3}\cline{5-7}
\multirow{2}{*}{$ C_9 $} & $ \ne 3 $ & $ 3 \cdot 9^3 $ & & \multirow{2}{*}{$ C_{14} $} & $ 2 $ & $ 2 \cdot 14^3 $ \\
\cline{2-3}\cline{6-7}
& $ 3 $ & $ 9^3 $ & & & $ 7 $ & $ 7^3 \cdot 14^3 $ \\
\cline{1-3}\cline{5-7}
\multirow{3}{*}{$ C_{10} $} & $ \ne 2, 5 $ & $ 5 \cdot 10^3 $ & & \multirow{2}{*}{$ C_{15} $} & $ 3 $ & $ 3 \cdot 15^4 $ \\
\cline{2-3}\cline{6-7}
& $ 2 $ & $ 2 \cdot 10^2 $ & & & $ 5 $ & $ 5^2 \cdot 15^4 $ \\
\cline{2-3}\cline{5-7}
& $ 5 $ & $ 5^2 \cdot 10^2 $ & & $ C_{18} $ & $ 2 $ & $ 2^2 \cdot 3 \cdot 18^3 $ \\
\cline{1-3}\cline{5-7}
\multirow{3}{*}{$ C_{12} $} & $ \ne 2, 3 $ & $ 12^4 $ & & $ C_{10} \oplus C_2 $ & $ 5 $ & $ 10^6 $ \\
\cline{2-3}\cline{5-7}
& $ 2 $ & $ 2^2 \cdot 12^2 $ & & $ C_{10} \oplus C_5 $ & $ 2 $ & $ 10^{12} $ \\
\cline{2-3}\cline{5-7}
& $ 3 $ & $ 3 \cdot 6 \cdot 12^2 $ & & $ C_{12} \oplus C_2 $ & $ 3 $ & $ 6^2 \cdot 12^4 $ \\
\cline{1-3}\cline{5-7}
\end{tabular}
\end{center}

\end{frame}

\begin{frame}{Proof strategy}

If $ P $ is a point of order $ N \ge 4 $, then the Tate normal form gives a morphism $ f : C \to X_1(N) $, which is surjective since $ j \notin k $.

\begin{example}[$ K = k(t) $ and $ \tor(E) \cong C_5 $]
The morphism $ f : \P_k^1 \to X_1(5) \cong \P_f^1 $ has a zero $ v_1 $ and a pole $ v_2 $, lying over the cusps $ 0 $ and $ \infty $ of $ X_1(5) $, which are distinct places of $ K $.

\bigskip As before, $ c_{v_1}(E) = v_1(\Delta) = 5v_1(f) $ at $ v_1 $.

\bigskip On the other hand, the change of variables $ (x, y) \mapsto (x / f^2, y / f^3) $ gives
$$ y^2 - (1 - f^{-1})xy - f^{-2}y = x^3 - f^{-1}x^2, \qquad \Delta = f^{-5} - 11f^{-6} - f^{-7}, $$
whose reduction at $ v_2 $ is $ y^2 - xy = x^3 $, which is also split multiplicative with $ c_{v_2}(E) = -5v_2(f) $.
\end{example}

\end{frame}

\begin{frame}{McDonald parameterisations}

\begin{theorem}[Cox--Parry 1980, McDonald 2018]
Let $ E $ be a non-isotrivial elliptic curve over $ k(t) $ with a point of order $ N \ge 4 $. Then $ \tor(E) \cong T $ with an explicit Tate normal form.
\begin{center}
\begin{tabular}{|r|l|}
\hline
$ \ch(k) $ & $ T $ \\
\hline
any & $ C_4 $, $ C_5 $, $ C_6 $, $ C_7 $, $ C_8 $, $ C_9 $, $ C_{10} $, $ C_{12} $ \\
\hline
$ \ne 2 $ & $ C_4 \oplus C_2 $, $ C_6 \oplus C_2 $, $ C_8 \oplus C_2 $, $ C_4 \oplus C_4 $ \\
\hline
$ \ne 3 $ & $ C_6 \oplus C_3 $ \\
\hline
$ \ne 5 $ & $ C_5 \oplus C_5 $ \\
\hline
$ 2 $ & $ C_{14} $, $ C_{18} $, $ C_{10} \oplus C_5 $ \\
\hline
$ 3 $ & $ C_{15} $, $ C_{12} \oplus C_2 $ \\
\hline
$ 5 $ & $ C_{15} $, $ C_{10} \oplus C_2 $ \\
\hline
$ 7 $ & $ C_{14} $ \\
\hline
$ 11 $ & $ C_{11} $ \\
\hline
\end{tabular}
\end{center}
\end{theorem}

\end{frame}

\section{Applications of AI}

\begin{frame}{Lean formalisation}

As part of Axiom Math, I manually formalised Tate's algorithm over perfect residue fields, which \emph{defines} the Tamagawa number.

\bigskip This definition works over any commutative ring equipped with a reduction datum and an additive valuation that are compatible, including $ \Z_p $ and $ k[[t]] $, but also $ \Z $ and $ k[t] $ upon choosing a uniformiser.

\bigskip This definition is good for proofs, but also for kernel computation given a strategy to compute $ p $-th roots and roots of cubics when $ \ch(k) = p $.

\bigskip Using GPT-5.6 Sol, I autoformalised the divisibility results and the sharp examples over a few weeks, with refactoring still ongoing with Fable 5.1.

\bigskip The total lines of code is approximately $ 3,000 $ lines for Tate's algorithm, $ 30,000 $ lines for divisibility results, and $ 30,000 $ lines for sharp examples.

\end{frame}

\begin{frame}{Reference errors}

GPT-6 Astra detected minor mistakes in the literature.

\bigskip Ulmer 2019 claims that $ \#\tor(E)^2 $ divides $ \Reg(E) \cdot c(E) $, based on a claimed proof that $ c(E) = D(E) $ by Berger et al 2020. However, the elliptic curve $ E $ over $ \F_3(t) $ given by $ y^2 + xy + (t^3 - t + 1)y = x^3 $ has
$$ \dfrac{\Reg(E) \cdot c(E)}{\#\tor(E)^2} = \dfrac{1}{3}. $$
McDonald 2018 claims that $ \tor(E) \cong C_{12} \oplus C_2 $ when $ \ch(k) = 3 $ can only occur when $ \zeta_4 \in k $. However, the elliptic curve $ E $ over $ \F_3(t) $ given by
$$ y^2 + xy + \dfrac{t^3(t^2 - 1)^3}{(t^2 + 1)^6}y = x^3 + \dfrac{t^3(t^2 - 1)^3}{(t^2 + 1)^6}x^2 $$
has $ \tor(E) \cong C_{12} \oplus C_2 $, but $ \zeta_4 \notin \F_3 $.

\end{frame}

\begin{frame}{Future work}

Observe that $ 5^2 \mid c(E) $ when $ \tor(E) \cong C_5 $, but also $ 7^3 \mid c(E) $ when $ \tor(E) \cong C_7 $ and $ 11^5 \mid c(E) $ when $ \tor(E) \cong C_{11} $.

\begin{conjecture}
A point of prime order $ N \ge 5 $ forces $ N^{(N - 1) / 2} \mid c(E) $.
\end{conjecture}

\begin{conjecture}
Independently of the Birch and Swinnerton-Dyer conjecture,
$$ \dfrac{q^{h + g - 1}}{\Reg(E)} \cdot L^*(E) \in \Z[\tfrac{1}{6}], \qquad \dfrac{c(E)}{\#\tor(E)^2} \in \Z[\tfrac{1}{6}]. $$
\end{conjecture}

GPT-6 Astra claims $ q^{h + g - 1} \cdot L^*(E) \in \Z $ by calculating Newton slopes of $ L(E, T) $ as the characteristic polynomial of Frobenius acting on the crystalline cohomology of the surface associated to $ E $.

\end{frame}

\end{document}